\clearpage
\section{Area probability with an angle}
\begin{colbox}
	A needle of length \(l\) is dropped on a grid of parallel lines with distance \(a\) between them, such that \(l<a\).
	\begin{enumerate}
		\item What is the chance the needle hits the net if it is dropped such that it must be perpendicular to the net?
		\item What is the chance it hits the net if it fells at a random angle?
	\end{enumerate}
\end{colbox}
\subsection{Perpendicular needle}
In this case it's simply the length of the needle divided by the size of the gap
\begin{equation}
	P=\frac{l}{a}
\end{equation}
\subsection{at angle}
In this case there actually needs to be a calculation. First we notice that given we already know the angle at which the needle dropped, the chance is simply
\begin{equation}
	P(hit) = \frac{l}{a}\sin\theta
\end{equation}
therefore
\begin{align}
	P(hit) &= \int_{0}^{\pi}\dd{\theta}P(hit|\Theta=\theta)P(\Theta=\theta)\\
	&= \frac{l}{a}\int_{0}^{\pi}\dd{\theta}\sin\theta P(\Theta=\theta)\\
	&= \frac{2l}{\pi a}
\end{align}